\chapter{The Sun \& Heliophysics}

\section*{Theory problems}

\probhead{4.1}{}{2008}{Bandung, Indonesia}{TH S2}{20}
% source id: 2008_TH_S2   tag: sunspot area in msh, foreshortening
On April 2, 2008 a telescope (10\,cm diameter, $f/10$) at the Bosscha
Observatory was used to observe the Sun and found an active region 0987 (based
on the NOAA number) at $8^\circ$ South and $40^\circ$ West from the centre of
the solar disk. The region was recorded with a CCD SBIG ST-8 Camera
($1600 \times 1200$ pixels, $(9\,\mu\mathrm{m} \times 9\,\mu\mathrm{m})$ per
pixel) and its size was $5 \times 4$ pixels. According to the Astronomical
Almanac, the solar diameter is $32'$. How large is the corrected area of the
active region in units of millionths of the solar hemisphere (msh)?

\probhead{4.2}{}{2010}{Beijing, China}{TH S13}{10}
% source id: 2010_TH_S13   tag: solar photosphere temperature from color
Estimate the effective temperature of the photosphere of the Sun using the naked
eye colour of the Sun.

\probhead{4.3}{}{2011}{Katowice, Poland}{TH S14}{10}
% source id: 2011_TH_S14   tag: solar neutrino flux from fusion
Estimate the number of solar neutrinos which should pass through a
$1\,\mathrm{m}^2$ area of the Earth's surface perpendicular to the Sun every
second. Use the fact that each fusion reaction in the Sun produces
$26.8$\,MeV of energy and 2 neutrinos.

\probhead{4.4}{}{2012}{Rio de Janeiro, Brazil}{TH S6}{}
% source id: 2012_TH_S6   tag: p-p chain He-3 fusion energy release
Most of the energy emitted by the Sun is generated in its core via the so-called
proton--proton (p--p) nuclear chain reaction, which has three different
branches. The most energetic branch transforms $2\,^{3}\mathrm{He}$ into
$^{4}\mathrm{He} + 2\,^{1}\mathrm{H}$. Calculate the energy released (in MeV)
and the fractional reduction of the mass of the particles involved in this
reaction.


\probhead{4.5}{Sunspot}{2018}{Beijing, China}{TH T7}{25}
% source id: 2018_TH_T7   tag: sunspot magnetic pressure equilibrium
Magnetic fields are important in the physics of stars and sunspots. As an
approximation, we can model the photosphere of the Sun as consisting of a
plasma, which can be simply treated as a single component ideal gas, and a
magnetic field ($B$), which has an associated magnetic pressure
\[ p_B = \frac{B^2}{2\mu_0}. \]
It behaves like any other physical pressure except that it is carried by the
magnetic field rather than by the kinetic energy of particles.

Assume that the number density of particles in the photosphere is constant
everywhere, but the magnetic field inside the sunspot
($B_{\mathrm{in}} = 0.1$\,T) is much stronger than outside
($B_{\mathrm{out}} = 5 \times 10^{-3}$\,T). From the blackbody spectrum, the
temperature inside the sunspot is $T_{\mathrm{in}} \sim 4000$\,K, while the
temperature outside is $T_{\mathrm{out}} \sim 6000$\,K (which is why the sunspot
looks darker). For the sunspot to be stable, the inside must be in equilibrium
with the outside.

\begin{description}
\item[(a)] Estimate the number density of plasma particles in the solar
  photosphere.
\item[(b)] Compare your answer with an estimate of the number density of
  particles in the atmosphere at the surface of the Earth.
\end{description}

\probhead{4.6}{Effect of sunspots on solar irradiance}{2019}{Keszthely, Hungary}{TH 7}{20}
% source id: 2019_TH_7   tag: solar constant with sunspot
Since 1978, the solar constant has been almost continuously measured by
detectors on-board artificial satellites. These accurate measurements revealed
that there are seasonal, monthly, yearly, and longer timescale variations in the
solar constant. While the seasonal variations have their origins in the
periodically varying Earth--Sun distance, the decade-long quasi-cyclic
variations mainly depend on the activity cycle(s) of the Sun.

\begin{description}
\item[(a)] Calculate the value of the solar constant at the top of the Earth's
  atmosphere, when the Earth is 1\,au from a perfectly quiet Sun, assuming that
  the Sun emits as a perfect black-body. \hfill[4\,p]
\item[(b)] Calculate the solar constant of this perfectly quiet Sun in early
  January and early July, and find their ratio. \hfill[4\,p]
\item[(c)] Calculate the solar constant again at 1\,au in the presence of a
  near-equatorial sunspot with mean temperature $T_{\mathrm{sp}} = 3300$\,K and
  diameter $D_{\mathrm{sp}} = 90\,000$\,km. Calculate the ratio --- blank Sun to
  Sun with sunspot. \hfill[7\,p]\\
  Assume the sunspot is circular and ignore the effects of its spherical
  projection. Neglect any other activity features. Also assume that the Sun is
  rotating fast enough, hence solar irradiance is still isotropic.
\item[(d)] In reality solar irradiance is no longer isotropic. Calculate the
  ratio of solar irradiance for the case when the sunspot is not visible from
  the Earth to the case when it is fully visible. \hfill[5\,p]
\end{description}

\probhead{4.7}{Under Pressure}{2021}{Bogot\'a (remote)}{TH Q5}{10}
% source id: 2021_TH_Q5   tag: solar flux-tube magnetic vs gas pressure
Magnetic fields in the Sun are constantly shaping the structure of various
different features in the solar atmosphere. Inside any feature, the magnetic
field ($B$) adds to the total pressure exerted by the gas. This so-called
magnetic pressure is a function of the height $z$ and can be expressed as
\[ P_{mag}(z) = \frac{B^2(z)}{2\mu_0}. \]
On the other hand, the gas can be considered to be in hydrostatic equilibrium
and hence the gas pressure decays exponentially from an initial pressure value
$P_0$ with increasing $z$. It can be expressed as
\[ P_{gas}(z) = P_0 e^{-z/H}, \]
where $H$ is the scale height, i.e.\ the height at which the pressure falls to
$P_0/e$.

Consider one type of feature, a magnetic flux tube rising from the solar surface
up into an unmagnetized environment (see figure below). Assuming that the total
pressure of the material inside the tube and of the material outside it is in
equilibrium:

\begin{description}
\item[(5.1)] Find an expression for the magnetic field strength as a function of
  height $z$. \hfill[7.0\,pt]
\item[(5.2)] If the magnetic field at the base of a flux tube is $0.3$\,T, and
  the scale height $H$ in a given solar model is $150$\,km, at what height will
  the magnetic field be reduced to $0.03$\,T? \hfill[3.0\,pt]
\end{description}

\ioaafig{2021_TH_Q5-f1.pdf}

\section*{Data-analysis problems}

\probhead{4.8}{}{2013}{Volos, Greece}{DA D2}{}
% source id: 2013_DA_D2   tag: sunspot tracking, solar synodic rotation
You are given 5 recent photographs of the solar photosphere shot at exactly the
same time every two days (May 1 -- May 9, 2013) in equatorial coordinates. You
are also given two transparent Stonyhurst grids, which display heliocentric
coordinates (heliocentric longitude $\ell$, and heliocentric latitude $b$). They
cover the interval between April 28 to May 15. As the Earth does not orbit
exactly around the Sun's equator, through the year the solar equator seems to
move up and down a little more than 7 degrees from the centre of the solar disc.
This angle, $B_0$, varies sinusoidally through the year. Furthermore, the axis
of rotation of the Sun, as seen from the Earth, does not coincide with the axis
of rotation of the Earth. The angle on the plane of the sky between the two
axes, $P_0$, also varies through the year. The numerical values of these angles
($B_0$ and $P_0$) are indicated on each of the 5 images of the Sun.

\begin{enumerate}
\item Mark the axis of rotation of the Sun on each photograph.
\item Choose 3 prominent sunspots that can be followed in all (or most)
  photographs and mark them as $S_1$, $S_2$ and $S_3$ on the photos. Using the
  appropriate Stonyhurst grids, find their coordinates $(\ell, b)$ for every day
  (May 1 to May 9) and note them down in Table 1.
\item Construct the diagrams $\Delta\ell/\Delta t$ for each sunspot.
\item Calculate its synodic period ($P$) of rotation \textbf{in days} for each
  sunspot. Write down the result for each sunspot, $P_{S1}$, $P_{S2}$,
  $P_{S3}$.
\item Calculate the average synodic period ($P_\odot$) of rotation of the Sun
  in days.
\end{enumerate}

\probhead{4.9}{Predicting arrival times of coronal mass ejections on Earth}{2025}{Mumbai, India}{DA D02}{60}
% source id: 2025_DA_D02   tag: CME drag model arrival prediction
The Sun occasionally releases magnetized plasma, termed coronal mass
ejections (CMEs), that originate from the surface of the Sun and propagate
outwards. Accurate prediction of their arrival times at Earth is crucial for
understanding and mitigating their potential effects on satellites orbiting
the Earth. In this problem, we aim to predict the arrival times of CMEs by
developing an empirical model, using the data of 10 CMEs. Throughout this
problem, the distance between the Sun's surface and Earth is taken to be
$214\,R_\odot$. Further, assume that the Sun is not rotating. Due to
electromagnetic, gravitational and drag forces, CMEs experience a variable
acceleration throughout their propagation. In the first two parts of this
problem, we assume that the region between the Sun and the Earth is vacuum.

\medskip\noindent\textbf{CMEs through vacuum}

\begin{description}
\item[(D02.1)] The initial velocity, $u$, at the solar surface
  ($= 1\,R_\odot$), the final velocity, $v$, upon reaching Earth, and the time
  to arrive at Earth after leaving the surface of the Sun (in hours), $\tau$,
  are given for 10 CMEs in the following table.

  \begin{center}
  \begin{tabular}{l c c c}
  \toprule
  CME & $u$ & $v$ & $\tau$\\
  Name & (km\,s$^{-1}$) & (km\,s$^{-1}$) & (h)\\
  \midrule
  CME-A & 804  & 470 & 74.5\\
  CME-B & 247  & 360 & 127.5\\
  CME-C & 523  & 396 & 103.5\\
  CME-D & 830  & 415 & 71.0\\
  CME-E & 665  & 400 & 104.5\\
  CME-F & 347  & 350 & 101.5\\
  CME-G & 446  & 375 & 99.5\\
  CME-H & 155  & 360 & 97.0\\
  CME-I & 1016 & 515 & 67.0\\
  CME-J & 683  & 410 & 54.0\\
  \bottomrule
  \end{tabular}
  \end{center}
\item[(D02.1a)] Calculate the average acceleration, $a$, for each CME in
  m\,s$^{-2}$. \hfill[3]
\item[(D02.1b)] We assume an empirical model for the acceleration,
  $a_{\mathrm{model}}$, of a CME, which depends on its initial velocity, $u$,
  as $a_{\mathrm{model}} = m\left(\frac{u}{u_0}\right) + \alpha$, where
  $a_{\mathrm{model}}$ is expressed in m\,s$^{-2}$, $u$ is expressed in
  km\,s$^{-1}$ and $u_0 = 1.00 \times 10^{3}$\,km\,s$^{-1}$. Determine the
  constants $m$ and $\alpha$ and their associated uncertainties using an
  appropriate graph. \hfill[15]
\item[(D02.1c)] For each CME, tabulate $a_{\mathrm{model}}$ in m\,s$^{-2}$.
  Hence calculate the root-mean-square (rms) deviation of acceleration,
  $\delta a_{\mathrm{rms}}$, between the calculated accelerations, $a$, and
  the model values, $a_{\mathrm{model}}$. \hfill[4]
\item[(D02.2)] We consider two other CMEs: CME-1 and CME-2, with initial
  velocities, $u = 1044$\,km\,s$^{-1}$ and $273$\,km\,s$^{-1}$, respectively.
\item[(D02.2a)] Using the empirical model obtained in (D02.1b), calculate the
  predicted arrival times at Earth, $\tau_{1,m}$ and $\tau_{2,m}$ (in hours),
  for CME-1 and CME-2, respectively. \hfill[4]
\item[(D02.2b)] The observed arrival times at Earth of CME-1 and CME-2 are
  46.0\,h and 74.5\,h, respectively. The empirical model is considered to be
  VALID for a particular CME if its predicted arrival time is within 20\% of
  its observed arrival time; otherwise, it is NOT VALID. Indicate the validity
  of the model for each CME. \hfill[2]
\end{description}

\medskip\noindent\textbf{CMEs in presence of solar wind}

In reality, the space between the Sun and the Earth is permeated with the
solar wind, which exerts a drag force on CMEs. This drag force can either
decelerate or accelerate a CME, depending on the CME's velocity relative to
that of the solar wind. To account for the solar wind's influence, we will use
a ``drag-only'' model for distances $R_{\mathrm{obs}}(t) \geq R_0$, where
$R_0$ is the distance beyond which the drag force becomes the dominant force
affecting the CME's motion.

The distance from the surface of the Sun as determined from the ``drag-only''
model, $R_D(t)$, and velocity, $V_D(t)$, of a CME in this model is given by
\begin{gather*}
R_D(t) = \frac{S}{\gamma}\ln\left[1 + S\gamma(V_0-V_s)(t-t_0)\right]
  + V_s(t-t_0) + R_0\\
V_D(t) = \frac{V_0 - V_s}{1 + S\gamma(V_0-V_s)(t-t_0)} + V_s
\end{gather*}
where $\gamma = 2\times10^{-8}$\,km$^{-1}$, $V_s$ is the constant speed of the
solar wind, $R_0$ and $V_0$ are the distance and velocity, respectively, at
time $t_0$, and $S$ is the sign factor: $S = 1$ if $V_0 > V_s$; $S = -1$ if
$V_0 \leq V_s$.

\begin{description}
\item[(D02.3)] The tables below show the observed radial distance from the
  surface of the Sun, $R_{\mathrm{obs}}(t)$ (measured in $R_\odot$), as a
  function of time, $t$ (in hours), for two CMEs: CME-3 and CME-4. The last
  data point in each table (D5 and P8, respectively) corresponds to the
  arrival time of the respective CME at Earth. For this part, assume
  $V_s = 330$\,km\,s$^{-1}$.

  \begin{center}
  \begin{tabular}[t]{c c c}
  \multicolumn{3}{c}{CME-3}\\
  \toprule
  Data point & $t$ (in h) & $R_{\mathrm{obs}}(t)$ (in $R_\odot$)\\
  \midrule
  D1 & 0.200 & 6.36\\
  D2 & 0.480 & 7.99\\
  D3 & 1.22  & 11.99\\
  D4 & 1.49  & 13.51\\
  D5 & 58.05 & 214\\
  \bottomrule
  \end{tabular}
  \quad
  \begin{tabular}[t]{c c c}
  \multicolumn{3}{c}{CME-4}\\
  \toprule
  Data point & $t$ (in h) & $R_{\mathrm{obs}}(t)$ (in $R_\odot$)\\
  \midrule
  P1 & 1.00 & 4.00\\
  P2 & 3.00 & 6.00\\
  P3 & 4.00 & 9.00\\
  P4 & 5.00 & 11.0\\
  P5 & 21.0 & 43.0\\
  P6 & 50.0 & 100\\
  P7 & 85.0 & 170\\
  P8 & 111  & 214\\
  \bottomrule
  \end{tabular}
  \end{center}

  We shall evaluate if the ``drag-only'' model satisfactorily predicts the
  arrival times of these CMEs. To use this model an appropriate choice of
  $t_0$, and corresponding $R_0$ and $V_0$ needs to be made.
\item[(D02.3a)] For CME-3, take the following two cases: (C1) $t_0$ is taken as
  the midpoint of the interval D1--D2; (C2) $t_0$ is taken as the midpoint of
  the interval D3--D4. Assume the velocity remains constant in each specific
  interval D1--D2 and D3--D4, but may differ between the two intervals. Using
  $t_0$, $R_0$, and $V_0$, calculate the difference between the observed and the
  predicted radial distance,
  $\delta R_D \equiv R_{\mathrm{obs}}(t) - R_D(t)$ in units of $R_\odot$ at
  $t = 58.05$\,h, for each of the two cases. \hfill[6]
\item[(D02.3b)] Evaluate $R_D(t)$ at points P5, P6, P7 and P8 between the Sun
  and the Earth for CME-4 for the following two cases, adopting a procedure
  similar to (D02.3a): (C3) $t_0$ is taken as the midpoint of the interval
  P1--P2; (C4) $t_0$ is taken as the midpoint of the interval P3--P4.
  \hfill[4]
\item[(D02.3c)] Plot $R_D(t)$ (in $R_\odot$) vs $t$ (in hours) for the two
  cases, C3 and C4, for CME-4 at points P5, P6, P7 and P8. On the same graph,
  draw smooth curves of $R_D(t)$ for the two cases. For this part, take the
  range of the $x$ axis from 0 to 180\,h. \hfill[10]
\item[(D02.3d)] Using the graph, estimate the absolute difference
  $|\delta\tau|$ between the actual arrival time of CME-4 at the Earth and its
  time of arrival predicted by the drag-only model, for each of the cases C3 and
  C4. \hfill[4]
\item[(D02.3e)] Indicate whether the following statement is TRUE or FALSE (no
  written justification needed): ``The drag forces exerted by the solar wind on
  CMEs become dominant for CME-3 at an earlier time compared to CME-4.''
  \hfill[1]
\item[(D02.4)] Consider drag as the dominant force acting on the 10 CMEs in part
  D02.1. Assume that the ``drag-only'' model is applicable from the surface of
  the Sun ($R_0 = 1\,R_\odot$) and beyond, for all CMEs. Estimate and tabulate
  the solar wind speed $V_s$ in km\,s$^{-1}$ for each CME. Further, estimate the
  average solar wind speed $V_{s,\mathrm{avg}}$ for all 10 CMEs. \hfill[7]
\end{description}

\section*{Observation --- written and sky map}

\probhead{4.10}{Solar Physics: the solar slap}{2021}{Bogot\'a (remote)}{OBS-SKY 1}{75}
% source id: 2021_OBS-SKY_1   tag: CME evolution from solar images
Using the provided series of solar images, select the appropriate images showing
the evolution of a coronal mass ejection (CME) and identify the CME front in
each. From these measurements, characterise the evolution of the CME.

\clearpage
\section*{Solutions to Chapter 4}

% Solutions are not transcribed. The official solutions for every theory and
% data-analysis problem are reproduced in full in the crop edition,
% ../IOAA_Problems_by_Topic_2007-2025.pdf, under the same problem number.

\solhead{4.1}{}
% source id: 2008_TH_S2
The pixel scale is
\begin{equation*}
\text{Pixel size}\,('') = \frac{206265 \times \text{pixel size}\,(\mu\mathrm{m})}{1000 \times \text{focal length (mm)}}
= \frac{206265 \times 9}{1000 \times 1000} = 1.86'' \tag*{(30\%)}
\end{equation*}
Then
\[
\text{Solar diameter} = 32 \times 60''/1.86''\,\text{per pixel} = 1034~\text{pixels} = 9306~\mu\mathrm{m}
\]
% sic: 32 x 60"/1.86" = 1032 pixels; the official solution states 1034 pixels
\begin{equation*}
\text{Solar radius,}\ R = 4653~\mu\mathrm{m} \tag*{(25\%)}
\end{equation*}
If $A_s$ is the measured area, then
\begin{equation*}
A_s = 5 \times 4 \times 81~\mu\mathrm{m}^2 = 1620~\mu\mathrm{m}^2 \tag*{(15\%)}
\end{equation*}
The corrected area due to the projection effect is
\begin{align*}
A_m &= A_s \cdot 10^6/(2\pi R^2 \cos B \cos L)\\
    &= 1620 \times 10^6/(2 \times 3.14 \times 4653^2 \times \cos 8^\circ \times \cos 40^\circ)\\
    &= 15.7~\mathrm{msh} \tag*{(30\%)}
\end{align*}

\solhead{4.2}{}
% source id: 2010_TH_S13
The Wien law is
\begin{equation*}
\lambda_{\max} = \frac{0.29}{T}\,(\mathrm{cm}) \tag*{5}
\end{equation*}
So the temperature is
\begin{equation*}
T = \frac{0.29}{550\times 10^{-9}} = 5272 \approx 5300~\mathrm{K} \tag*{5}
\end{equation*}
% sic: the official solution states the Wien constant in cm K but writes the
% wavelength as 550e-9 (and 500e-9 below), i.e. in metres; the quoted results
% 5272 K and 5800 K correspond to the wavelength expressed in cm.
Or
\[
T = \frac{0.29}{500\times 10^{-9}} = 5800~\mathrm{K}
\]
Note: $5200\sim6000$\,K all full mark.

\solhead{4.3}{}
% source id: 2011_TH_S14
The solar constant (from the table) is $1366$\,W\,m$^{-2}$.

Each reaction gives
$26.8\,\mathrm{MeV} = 26.8 \cdot 1.602 \cdot 10^{-13}\,\mathrm{J}
= 4.3 \cdot 10^{-12}\,\mathrm{J}$. \hfill[4]

Thus there must be $1366 / 4.3 \cdot 10^{-12} = 3.2 \cdot 10^{14}$ reactions per
second to produce this energy flux. \hfill[4]

Each reaction produces 2 neutrinos, so there must be
$6.4 \cdot 10^{14}$ neutrinos m$^{-2}\,$s$^{-1}$. \hfill[2]

\solhead{4.4}{}
% source id: 2012_TH_S6
The reaction chain is
\[
2\,\mathrm{He}^{3} \rightarrow \mathrm{He}^{4} + 2\,\mathrm{H}^{1} + E
\]
So the energy is
\begin{equation*}
E = \left(2m_{\mathrm{He}^{3}} - m_{\mathrm{He}^{4}} - 2m_{\mathrm{H}^{1}}\right)\cdot c^{2}
= 12.66~\mathrm{MeV} \tag*{(7 pt)}
\end{equation*}
\begin{equation*}
\frac{\Delta m}{m} = \frac{12.66~\mathrm{MeV}/c^{2}}{2m_{\mathrm{He}^{3}}} = 0.002254 \tag*{(3 pt)}
\end{equation*}

\solhead{4.5}{Sunspot}
% source id: 2018_TH_T7
\begin{description}
\item[(a)] The pressure is from all directions. The plasma inside the sunspot
  must have the same total pressure as the plasma outside to maintain
  equilibrium on the border.

  The kinetic pressure inside the sunspot is $p_i = n_i k_B T_i$, where $n_i$
  is the number density inside the sunspot. And also,
  $p_e = n_e k_B T_e$. \hfill(3 points)

  ($p_e V = N_e k_B T_e$ is also correct. But it's better to write the
  pressure in terms of number density.)

  From the assumption, $n_i = n_e = n$. The equilibrium
  requires: \hfill(3 points)
  \begin{equation*}
  n k_B T_i + \frac{B_i^2}{2\mu_0} = n k_B T_e + \frac{B_e^2}{2\mu_0}
  \tag*{(8 points)}
  \end{equation*}
  Then
  \begin{equation*}
  n = \frac{1}{2\mu_0 k_B}\,\frac{(B_i^2 - B_e^2)}{(T_e - T_i)}
    = 1.43 \times 10^{23}\,m^{-3}
  \tag*{(6 points)}
  \end{equation*}
\item[(b)] On the Earth, $p_E = n_E k_B T_E$ where $P_E \approx 10^5$\,Pa and
  $T_E \approx 300$\,K.

  Thus,
  \begin{equation*}
  n_E = \frac{p_E}{k_B T_E} = 2.4 \times 10^{25}\,m^{-3}
  \tag*{(4 points)}
  \end{equation*}
  This means that the number density of the atmospheric particles at the
  surface of the earth is at least 100 times larger that the number density of
  particles in the solar photosphere (of course the mass density is even
  higher). \hfill(1 point)
  % sic: "larger that" as in the original
\end{description}

\solhead{4.6}{Effect of sunspots on solar irradiance}
% source id: 2019_TH_7
\begin{description}
\item[(a)] The Sun's total power can be calculated from the
  Stefan--Boltzmann law: \hfill(1 p)
  \begin{equation*}
  P = A\sigma T^4, \tag{7.1}
  \end{equation*}
  where $A$ is the emitting surface, $\sigma$ is the Stefan--Boltzmann
  constant and $T$ is the (effective) temperature.

  The total power emitted homogeneously and isotropically by the Sun
  ($P_\odot$) distributes uniformly on the surface area of a sphere with
  radius $d_\oplus = 1$\,au. Hence: \hfill(2 p)
  \begin{equation*}
  I_\odot = \frac{P_\odot}{4\pi d_\oplus^2}
          = \frac{\sigma 4\pi R_\odot^2 T_\odot^4}{4\pi d_\oplus^2}
          = \sigma\,\frac{R_\odot^2 T_\odot^4}{d_\oplus^2}, \tag{7.2}
  \end{equation*}
  where $R_\odot$ and $T_\odot$ denote the radius and effective temperature of
  the Sun, respectively.

  With the numerical values taken from the Table of constants:
  $I_\odot = 1365.9$\,W\,m$^{-2}$. \hfill(1 p)
\item[(b)] In the first days of January the Earth is in perihelion, while in
  the beginning of July it is in aphelion. Therefore, instead of the mean
  Earth--Sun distance (i.e.\ the semi-major axis of the Earth's orbit), now
  one should substitute the perihelion and aphelion distances, respectively.
  Therefore, one has to write that: \hfill(1 p)
  \begin{gather*}
  I_{\odot,\mathrm{perihelion}} = \frac{P_\odot}{4\pi(d_\oplus(1-e))^2}
    = \sigma \frac{R_\odot^2 T_\odot^4}{(d_\oplus(1-e))^2}
    = \frac{I_\odot}{(1-e)^2}, \tag{7.3}\\
  I_{\odot,\mathrm{aphelion}} = \frac{P_\odot}{4\pi(d_\oplus(1+e))^2}
    = \sigma \frac{R_\odot^2 T_\odot^4}{(d_\oplus(1+e))^2}
    = \frac{I_\odot}{(1+e)^2}, \tag{7.4}
  \end{gather*}
  where $e = 0.016\,710\,22$ is the eccentricity of the Earth's orbit (see the
  Table of constants).

  With numerical values:

  in early January: $I_{\odot,\mathrm{perihelion}} = 1412.7$\,W\,m$^{-2}$
  \hfill(1 p)

  in early July: $I_{\odot,\mathrm{aphelion}} = 1321.4$\,W\,m$^{-2}$
  \hfill(1 p)

  The ratio of the values:
  $I_{\odot,\mathrm{perihelion}}/I_{\odot,\mathrm{aphelion}} = 1.0691$. The
  perihelion irradiation is higher by about 7\%. \hfill(1 p)
\item[(c)] In this case the total power of the Sun becomes \hfill(2 p)
  \begin{equation*}
  P'_\odot = \sigma\left[T_\odot^4(A_\odot - A_{\mathrm{sp}})
    + T_{\mathrm{sp}}^4 A_{\mathrm{sp}}\right], \tag{7.5}
  \end{equation*}
  where $A_\odot$ stands for the solar surface, while $A_{\mathrm{sp}}$ is the
  area of the spot (as a circle).

  Thus, the solar constant takes the following form: \hfill(2 p)
  \begin{align*}
  I'_\odot &= \frac{P'_\odot}{4\pi d_\oplus^2}
    = \frac{\sigma\pi\left[T_\odot^4(4R_\odot^2 - R_{\mathrm{sp}}^2)
      + T_{\mathrm{sp}}^4 R_{\mathrm{sp}}^2\right]}{4\pi d_\oplus^2}
    = \frac{\sigma}{4 d_\oplus^2}\left[T_\odot^4(4R_\odot^2 - R_{\mathrm{sp}}^2)
      + T_{\mathrm{sp}}^4 R_{\mathrm{sp}}^2\right] = \tag{7.6}\\
  &= I_\odot\left[1 - \left(\frac{R_{\mathrm{sp}}}{2R_\odot}\right)^2
    + \left(\frac{R_{\mathrm{sp}}}{2R_\odot}\right)^2
      \left(\frac{T_{\mathrm{sp}}}{T_\odot}\right)^4\right]
   = I_\odot\left[1 - \left(\frac{R_{\mathrm{sp}}}{2R_\odot}\right)^2
      \left(1 - \left(\frac{T_{\mathrm{sp}}}{T_\odot}\right)^4\right)\right] =\\
  &= 0.9991\,I_\odot \tag*{(2 p)}
  \end{align*}
  With numerical value: $I'_\odot = 1364.6$\,W\,m$^{-2}$. \hfill(1 p)
\item[(d)] As we have seen, the ratio is
  \[
  A = \frac{I'_\odot}{I_\odot}
    = \frac{1364.64~\mathrm{W\,m^{-2}}}{1365.92~\mathrm{W\,m^{-2}}} = 0.9991,
  \]
  or, converting it into percentage, the solar constant in the case of the
  spotted Sun will be lower by 0.09\%.

  When this sunspot emerges in the surface of the Sun, solar emission no
  longer will be isotropic, i.e.\ uniform in every directions. When the
  sunspot is fully seen from the Earth, the ratio of the irradiation of the
  Earth in the spotted and unspotted cases will be equal to the ratio of the
  solar powers in the two cases, i.e.,
  \begin{gather*}
  P_{\odot,\mathrm{disk}} = \sigma T_\odot^4 A_\odot
    = \sigma\pi T_\odot^4 R_\odot^2 \tag{7.7}\\
  P'_{\odot,\mathrm{disk}} = \sigma\left[T_\odot^4(A_\odot - A_{\mathrm{sp}})
    + T_{\mathrm{sp}}^4 A_{\mathrm{sp}}\right]
    = \sigma\pi\left[T_\odot^4(R_\odot^2 - R_{\mathrm{sp}}^2)
    + T_{\mathrm{sp}}^4 R_{\mathrm{sp}}^2\right] \tag{7.8}
  \end{gather*}
  \hfill(1 p each)

  and therefore, the ratio of the irradiations (or the effective solar
  constants) can be written as: \hfill(2 p)
  \begin{equation*}
  A' = \frac{P'_{\odot,\mathrm{disk}}}{P_{\odot,\mathrm{disk}}}
     = \frac{T_\odot^4(R_\odot^2 - R_{\mathrm{sp}}^2)
       + T_{\mathrm{sp}}^4 R_{\mathrm{sp}}^2}{T_\odot^4 R_\odot^2}
     = 1 - \left(\frac{R_{\mathrm{sp}}}{R_\odot}\right)^2
       \left[1 - \left(\frac{T_{\mathrm{sp}}}{T_\odot}\right)^4\right] \tag{7.9}
  \end{equation*}
  With numerical values: $A' = 0.9963$. \hfill(1 p)

  Thus, when the spot directed towards the Earth the total irradiation of the
  Earth becomes weaker by 0.37\%.
  % sic: "when the spot directed towards" as in the original
\end{description}

\solhead{4.7}{Under Pressure}
% source id: 2021_TH_Q5
\begin{description}
\item[(5.1)] Pressure inside the flux tube $=$ Pressure outside the flux tube
  (surroundings):
  \[
  (\text{Magnetic pressure} + \text{Gas pressure inside}) = \text{Gas pressure outside}
  \]
  Particularly,
  \begin{equation*}
  \frac{B_0^2}{2\mu_0} + P_{gas_{in}} = P_{gas_{out}} \tag*{[2 points]}
  \end{equation*}
  \begin{equation*}
  \frac{B_0^2}{2\mu_0} = P_{0_{out}} - P_{0_{in}} \tag*{Eq.~(1)}
  \end{equation*}
  Similarly, for a particular height $z$,
  \begin{equation*}
  \frac{B(z)^2}{2\mu_0} = e^{-z/H}\left(P_{0_{out}} - P_{0_{in}}\right) \tag*{Eq.~(2)\quad[2 points]}
  \end{equation*}
  Dividing equation (2) to (1),
  \[
  \frac{B(z)^2}{B_0^2} = e^{-z/H}
  \]
  To finally get
  \begin{equation*}
  B(z) = B_0 e^{-z/2H} \tag*{[3 points]}
  \end{equation*}
\item[(5.2)] At $z = H$, $B = B_0 e^{-z/150}$
  % sic: the official solution writes "At z = H, B = B_0 e^{-z/150}"; with
  % H = 150 km the exponent should read -z/(2 x 150) = -z/300, as used below.
  \begin{align*}
  0.03 &= 0.3\, e^{-\frac{z}{2\times150}}\\
  z &= 300 \times \ln 10\\
  z &= 300 \times 2.301\\
  % sic: ln 10 = 2.303; the official solution uses 2.301
  z &= 690~km \tag*{[3 points]}
  \end{align*}
\end{description}

\solhead{4.8}{}
% source id: 2013_DA_D2
Five large photographs of the solar photosphere (adapted from
www.spaceweather.com) and the appropriate (and on-scale) Stonyhurst grids will
be given to the students (Figure 3A). The photographs given to the students,
will not be annotated. Each photograph is accompanied by the angle $B_0$ and
$P_0$ of the day of the observation. The student should calculate the rotation
of the Sun by measuring the coordinates of at least 3 well recognized sunspots
as they follow the rotation of the Sun.

\ioaafig{2013_DA_D2-s1.pdf}

\begin{enumerate}
\item Correctly draw the axis of rotation of the Sun. This can be done by
  drawing a straight line at an angle of $P_0$ degrees anticlockwise from the
  vertical for each photograph (note: the photographs are given in equatorial
  coordinates.) \hfill(4 Points)
\item Choose the correct Stonyhurst grid and place it on each photograph, so
  that the solar axis on the grid coincides with the solar axis of the
  photograph. Estimate the heliocentric longitude, ($\ell_\odot$), and
  heliocentric latitude ($b_\odot$) with the help of the grid. Write down
  these coordinates in Table 1 for each of the 5
  photographs. \hfill(12 Points)

  \begin{center}
  \begin{tabular}{l cc cc cc}
  \toprule
  Date & \multicolumn{2}{c}{Sunspot S1}
    & \multicolumn{2}{c}{Sunspot S2} & \multicolumn{2}{c}{Sunspot S3}\\
  & $\ell_\odot$ & $b_\odot$ & $\ell_\odot$ & $b_\odot$
    & $\ell_\odot$ & $b_\odot$\\
  \midrule
  May 1 & 54    & 17.5 & 33  & 18.5 & 30.5 & 16\\
  May 3 & 27    & 18   & 6   & 19   & 4    & 16\\
  May 5 & 3     & 18.5 & $-19$ & 18.5 & $-23$ & 17\\
  May 7 & $-24$ & 18   & $-46$ & 18   & $-49$ & 16\\
  May 9 & $-49.5$ & 19 & $-72$ & 18   & $-77$ & 17\\
  \bottomrule
  \end{tabular}
  \end{center}
\item Construct the diagram $\Delta\ell_\odot/\Delta t$ for each sunspot.
  \hfill(2 for each graph)\quad(6 Points)

  \ioaafig{2013_DA_D2-s2.pdf}
\item Calculate the rotation ($\Delta\ell_\odot/\Delta t$) of the Sun for each
  sunspot:
  \[
  P_i = 8 * \frac{360}{\Delta\ell_{\odot i}}, \qquad
  P_1 = 27.8~days, \qquad
  P_2 = 27.4~days, \qquad
  P_3 = 26.8~days
  \]
  \hfill(3 for each graph)\quad(9 Points)
\item Calculate the average period of the Sun
  \begin{equation*}
  P_\odot = (P_1 + P_2 + P_3)/3 = 27.3~days \tag*{(1.5 Point)}
  \end{equation*}
\end{enumerate}

\solhead{4.9}{Predicting arrival times of coronal mass ejections on Earth}
% source id: 2025_DA_D02
\begin{description}
\item[(D02.1a)] Average acceleration (in m\,s$^{-2}$) is given by the
  following relation,
  \[
  a = (v-u)\times 10^{3}/(\tau\times 3600)
  \]
  \begin{center}
  \begin{tabular}{l c c c c}
  \toprule
  CME & $u$ & $v$ & $\tau$ & $a$\\
  Name & (km\,s$^{-1}$) & (km\,s$^{-1}$) & (h) & (m\,s$^{-2}$)\\
  \midrule
  CME-A & 804  & 470 & 74.5  & $-1.25$\\
  CME-B & 247  & 360 & 127.5 & 0.246\\
  CME-C & 523  & 396 & 103.5 & $-0.341$\\
  CME-D & 830  & 415 & 71.0  & $-1.62$\\
  CME-E & 665  & 400 & 104.5 & $-0.704$\\
  CME-F & 347  & 350 & 101.5 & 0.00821\\
  CME-G & 446  & 375 & 99.5  & $-0.198$\\
  CME-H & 155  & 360 & 97.0  & 0.587\\
  CME-I & 1016 & 515 & 67.0  & $-2.08$\\
  CME-J & 683  & 410 & 54.0  & $-1.40$\\
  \bottomrule
  \end{tabular}
  \end{center}
\item[(D02.1b)]
  \begin{align*}
  m &= (-3.05 \pm 0.09)~\mathrm{m\,s^{-2}}\\
  \alpha &= (1.08 \pm 0.06)~\mathrm{m\,s^{-2}}
  \end{align*}
  Thus, the expression obtained by fitting a straight line to $a$ vs $u/u_0$
  is
  \[
  a_{\mathrm{model}} = -3.05\left(\frac{u}{u_0}\right) + 1.08
  \]
  \textit{[the official solution shows a scanned hand-drawn fit on graph
  paper; not reproduced here]}
\item[(D02.1c)]
  \begin{center}
  \begin{tabular}{l c}
  \toprule
  CME & $a_{\mathrm{model}}$\\
  Name & (m\,s$^{-2}$)\\
  \midrule
  CME-A & $-1.37$\\
  CME-B & 0.327\\
  CME-C & $-0.515$\\
  CME-D & $-1.45$\\
  CME-E & $-0.948$\\
  CME-F & 0.0217\\
  CME-G & $-0.280$\\
  CME-H & 0.607\\
  CME-I & $-2.02$\\
  CME-J & $-1.00$\\
  \bottomrule
  \end{tabular}
  \end{center}
  The root mean square (rms) deviation of acceleration is given by
  \[
  \delta a_{\mathrm{rms}}
    = \sqrt{\frac{\sum_{i=1}^{10}(a_{\mathrm{model},i}-a_i)^2}{10}}
  \]
  Using this expression we get, $\delta a_{\mathrm{rms}} = 0.176$\,m\,s$^{-2}$.
\item[(D02.2a)] The time of arrival of a CME at the Earth, $\tau_m$, can be
  calculated from the following relation.
  \[
  s = u\tau_m + \frac{1}{2}a_{\mathrm{model}}\tau_m^2
  \]
  However, this would yield two solutions for $\tau_m$, only one of which will
  be for a positive final velocity (acceptable), and the other for a negative
  final velocity (not acceptable). Since acceleration is constant, the final
  velocity $v_m$ at Earth at a distance $s = 214\,R_\odot$ is given by
  \[
  v_m = \left(u^2 + 2a_{\mathrm{model}}s\right)^{1/2}
      = \left[u^2 + 2\left(m\frac{u}{u_0}+\alpha\right)s\right]^{1/2}
  \]
  Using this velocity, we get the arrival time $\tau$ as
  \[
  \tau_m = \frac{v_m - u}{a_{\mathrm{model}}}
         = \frac{v_m - u}{m\frac{u}{u_0}+\alpha}
  \]
  Using the given value of $u$, and the obtained values of $m$ and $\alpha$,
  we get

  For CME-1: $v_{1,m} = 680.728$\,km\,s$^{-1}$ which gives
  $\tau_{1,m} = 48.0$\,h.

  For CME-2: $v_{1,m} = 384.941$\,km\,s$^{-1}$ which gives
  $\tau_{2,m} = 126$\,h.
  % sic: the official solution writes v_{1,m} for CME-2 as well (should be
  % v_{2,m})
\item[(D02.2b)] The arrival time for CME-1 is $\tau_{1,m} = 48.0$\,h,
  therefore the percentage error in the arrival time of CME-1 is,
  \[
  \Delta\tau_{1,m} = \frac{48.0-46.0}{46.0}\times 100\% = 4.35\%
  \]
  Similarly, the arrival time for CME-2 is $\tau_{2,m} = 126$\,h, hence the
  percentage error is arrival time for CME-2 is,
  % sic: "error is arrival time" as in the original
  \[
  \Delta\tau_{2,m} = \frac{126-74.5}{74.5}\times 100\% = 69.1\%
  \]
  Hence CME-1: VALID; CME-2: NOT VALID.
\item[(D02.3a)] Inserting $R_0$, $t_0$, $V_0$, and $V_s$ in the given
  expression
  \[
  R_D(t) = \frac{S}{\gamma}\ln\left[1 + S\gamma(V_0-V_s)(t-t_0)\right]
         + V_s(t-t_0) + R_0,
  \]
  we get

  Case C1:
  \begin{align*}
  R_0 &= \frac{7.99+6.36}{2} = 7.18~R_\odot\\
  t_0 &= \frac{0.20+0.48}{2} = 0.340~\mathrm{h}\\
  V_0 &= \frac{(7.99-6.36)\times 695700}{(0.48-0.20)\times 3600}
       = 1.12\times 10^{3}~\mathrm{km\,s^{-1}}
  \end{align*}
  we get, $R_D = 210.61\,R_\odot$ at $t = 58.05$\,h. Thus,
  $\delta R_D = 214 - 210.61 = 3.39\,R_\odot$.

  Case C2:
  \begin{align*}
  R_0 &= \frac{11.99+13.51}{2} = 12.8~R_\odot\\
  t_0 &= \frac{1.22+1.49}{2} = 1.36~\mathrm{h}\\
  V_0 &= \frac{(13.51-11.99)\times 695700}{(1.49-1.22)\times 3600}
       = 1.09\times 10^{3}~\mathrm{km\,s^{-1}}
  \end{align*}
  Using the expression above, $R_D = 210.86\,R_\odot$ at $t=58.05$. Thus,
  $\delta R_D = 214 - 210.86 = 3.14\,R_\odot$.
\item[(D02.3b)] Case C3:
  \begin{align*}
  R_0 &= \frac{6+4}{2} = 5~R_\odot\\
  t_0 &= \frac{3+1}{2} = 2~\mathrm{h}\\
  V_0 &= \frac{(6-4)\times 695700}{(3-1)\times 3600} = 193~\mathrm{km\,s^{-1}}
  \end{align*}
  Case C4:
  \begin{align*}
  R_0 &= \frac{11+9}{2} = 10~R_\odot\\
  t_0 &= \frac{5+4}{2} = 4.5~\mathrm{h}\\
  V_0 &= \frac{(11-9)\times 695700}{(5-4)\times 3600} = 386~\mathrm{km\,s^{-1}}
  \end{align*}
  \begin{center}
  \begin{tabular}{c c c c}
  \toprule
  Data Point & $t$ (in h) & C3 & C4\\
  & & $R_D(t)$ (in $R_\odot$) & $R_D(t)$ (in $R_\odot$)\\
  \midrule
  P5 & 21.0 & 25.1 & 42.8\\
  P6 & 50.0 & 59.1 & 99.9\\
  P7 & 85.0 & 104  & 168\\
  P8 & 111  & 139  & 218\\
  \bottomrule
  \end{tabular}
  \end{center}
\item[(D02.3c)] The plot below shows the data points of CME-4 and smooth
  variation of $R_D$ extrapolated up to $220\,R_\odot$.

  \ioaafig{2025_DA_D02-s1.pdf}
\item[(D02.3d)] The arrival time of CME at Earth ($214\,R_\odot$) can be
  obtained by estimating the point of intersection between
  $R_D = 214\,R_\odot = $ constant and two smooth curves.

  The points of intersection are 167\,h and 109\,h for the cases C3 and C4,
  respectively. Thus, absolute time differences are:
  \[
  \text{C3: } |\delta\tau| = 56~\mathrm{h}
  \qquad\text{and}\qquad
  \text{C4: } |\delta\tau| = 2~\mathrm{h}.
  \]
\item[(D02.3e)] The statement is TRUE.

  In the scenario of CME-3, the ``drag-only'' model predicts the arrival
  distance accurately when $R_D$ is computed using the initial velocity,
  $V_0$ and initial distance, $R_0$ at 0.20\,h and beyond. While for CME-4,
  the ``drag-only'' model start accurately predicting the arrival times of
  CME-4 when $R_D$ is calculated using $V_0$ and $R_0$ at 4\,h after the
  launch of the CME.
  % sic: "model start" as in the original
\item[(D02.4)]
  \[
  V_D(t) = \frac{V_0 - V_s}{1 + S\gamma(V_0-V_s)(t-t_0)} + V_s
  \qquad\Longrightarrow\qquad
  v = \frac{u - V_s}{1 + S\gamma(u-V_s)\tau} + V_s.
  \]
  Taking, $S = 1$ if $v < u$, else $S = -1$. The above expression is
  quadratic equation, thus, we get two values of $V_s$. If $v < u$, choose
  the solution with $V_s < u$. If $v > u$, choose the solution with
  $V_s > u$.

  Calculating $V_s$ using the above conditions for all ten CMEs.
  \begin{center}
  \begin{tabular}{l c c c c}
  \toprule
  CME & $u$ & $v$ & $\tau$ & $V_s$\\
  Name & (km\,s$^{-1}$) & (km\,s$^{-1}$) & (h) & (km\,s$^{-1}$)\\
  \midrule
  CME-A & 804  & 470 & 74.5  & 337\\
  CME-B & 247  & 360 & 127.5 & 428\\
  CME-C & 523  & 396 & 103.5 & 314\\
  CME-D & 830  & 415 & 71.0  & 270\\
  CME-E & 665  & 400 & 104.5 & 303\\
  CME-F & 347  & 350 & 101.5 & 369\\
  CME-G & 446  & 375 & 99.5  & 305\\
  CME-H & 155  & 360 & 97.0  & 457\\
  CME-I & 1016 & 515 & 67.0  & 357\\
  CME-J & 683  & 410 & 54.0  & 248\\
  \bottomrule
  \end{tabular}
  \end{center}
  Taking the average of the solar wind speed, $V_s$, for each case, we get
  $V_{s,\mathrm{avg}} = 339$\,km\,s$^{-1}$.
\end{description}

\solhead{4.10}{Solar Physics: the solar slap}
% source id: 2021_OBS-SKY_1
\begin{enumerate}
\item \textbf{Data selection} (selecting appropriate images showing the
  evolution of the CME) and identification of CME front.

  \begin{center}
  \includegraphics[width=0.48\linewidth]{2021_OBS-SKY_1-s1-panel1.pdf}\hfill
  \includegraphics[width=0.48\linewidth]{2021_OBS-SKY_1-s1-panel2.pdf}\par\smallskip
  \includegraphics[width=0.48\linewidth]{2021_OBS-SKY_1-s1-panel3.pdf}\hfill
  \includegraphics[width=0.48\linewidth]{2021_OBS-SKY_1-s1-panel4.pdf}\par\smallskip
  \includegraphics[width=0.48\linewidth]{2021_OBS-SKY_1-s1-panel5.pdf}\hfill
  \includegraphics[width=0.48\linewidth]{2021_OBS-SKY_1-s1-panel6.pdf}\par\smallskip
  \includegraphics[width=0.48\linewidth]{2021_OBS-SKY_1-s1-panel7.pdf}
  \end{center}

  A total of 7 time intervals are selected tracing the evolution of the CME:
  \begin{enumerate}
  \item 2021/05/28 22:19:00 UT
  \item 2021/05/28 22:54:11 UT
  \item 2021/05/28 23:30:08 UT
  \item 2021/05/29 23:42:08 UT
  % sic: the official solution lists 2021/05/29 23:42:08 UT here, while its
  % summary table gives 2021/05/28 23:42:08 UT for interval 4
  \item 2021/05/29 01:12:16 UT
  \item 2021/05/29 03:36:08 UT
  \item 2021/05/29 05:00:08 UT
  % sic: the summary table gives 2021/05/29 05:36:08 UT for interval 7, but
  % 05:00:08 UT is consistent with the quoted 401 minutes from onset
  \end{enumerate}

  Summary of the information extracted from the images for all different
  selected intervals:

  {\footnotesize\setlength{\tabcolsep}{2.5pt}
  \begin{center}
  \begin{tabular}{c l c c c c c c}
  \toprule
  Interval & Date and Time & \begin{tabular}{@{}c@{}}Time\\ interval\\ (min)\end{tabular}
    & \begin{tabular}{@{}c@{}}Time from\\ CME onset\\ (min)\end{tabular}
    & \multicolumn{2}{c}{\begin{tabular}{@{}c@{}}Coordinates of\\ the CME front\end{tabular}}
    & \begin{tabular}{@{}c@{}}Distance from\\ disc center\\ (arcsec)\end{tabular}
    & \begin{tabular}{@{}c@{}}Distance from\\ disc center\\ (km)\end{tabular}\\
  & & & & $x''$ & $y''$ & &\\
  \midrule
  1 & 2021/05/28 22:19:00 UT & 0   & 0   & 950   & 144     & 961   & 696617\\
  2 & 2021/05/28 22:54:11 UT & 35  & 35  & 2682  & 174     & 2688  & 1948538\\
  3 & 2021/05/28 23:30:08 UT & 36  & 71  & 4579  & 97      & 4580  & 3320520\\
  4 & 2021/05/28 23:42:08 UT & 47  & 83  & 7441  & 116     & 7442  & 5395380\\
  5 & 2021/05/29 01:12:16 UT & 126 & 173 & 14000 & $-891$  & 14028 & 10170535\\
  6 & 2021/05/29 03:36:08 UT & 191 & 317 & 21048 & $-4570$ & 21538 & 15615349\\
  7 & 2021/05/29 05:36:08 UT & 210 & 401 & 27922 & $-7126$ & 28817 & 20892306\\
  \bottomrule
  \end{tabular}
  \end{center}
  % sic: the "Time interval" column of the official table is not consistent
  % with the differences of the successive timestamps (rows 4--7)

  \begin{center}
  \begin{tabular}{c c c c c c c}
  \toprule
  Interval & \begin{tabular}{@{}c@{}}Distance\\ from onset\\ (km)\end{tabular}
    & \begin{tabular}{@{}c@{}}Cumulative\\ velocity, plane\\ of the sky (km/s)\end{tabular}
    & \begin{tabular}{@{}c@{}}Deprojection of\\ velocity towards\\ ParkerSP (km/s)\end{tabular}
    & \begin{tabular}{@{}c@{}}Distance per\\ time interval\\ (km)\end{tabular}
    & \begin{tabular}{@{}c@{}}Velocity on\\ plane of the\\ sky (km/s)\end{tabular}
    & \begin{tabular}{@{}c@{}}Deprojection of\\ velocity towards\\ ParkerSP (km/s)\end{tabular}\\
  \midrule
  1 & 0        & 0   & 0    & 0       & 0   & 0\\
  2 & 1252198  & 596 & 728  & 1252198 & 596 & 728\\
  3 & 2624180  & 616 & 752  & 1371982 & 635 & 775\\
  4 & 4699040  & 944 & 1152 & 2074861 & 736 & 898\\
  5 & 9474195  & 913 & 1114 & 4775155 & 632 & 771\\
  6 & 14919009 & 784 & 958  & 5444814 & 475 & 580\\
  7 & 20195966 & 839 & 1025 & 5276957 & 419 & 511\\
  \bottomrule
  \end{tabular}
  \end{center}
  }
\item \textbf{Plotting the kinetic evolution of the CME.}

  Plots: distance vs time and velocity vs time.

  \ioaafig{2021_OBS-SKY_1-s2.pdf}
\item \textbf{Calculating the velocity of the CME front reaching ParkerSP.}

  The velocity corresponding to the last interval (up to reaching 30 solar
  radii) is 511\,km/s which, according to the considerations of the problem,
  would be the same velocity of the CME front once it reaches the ParkerSP.

  \textbf{Calculating the time spent by the CME front to reach the ParkerSP
  from its onset.}

  The Sun--ParkerSP distance estimated from Figure 1 is approx.\
  $0.66\,\mathrm{AU} = 9.9\times10^{7}\,\mathrm{km} = 141$ solar radii (given
  the information that 10 solar radii is approx.\ $7\times10^{6}$\,km).

  From the previous analysis, it is known that the CME front reaches a
  distance of 30 solar radii after 401 minutes (6.7 hours) from its onset
  (see Table). The distance that the CME travels at the constant speed of
  511\,km/s is expressed as $141-30=111$ solar radii
  $= 7.8\times 10^{7}$\,km, and thus the time spent in that interval is
  42.4 hours.

  Finally, the total time it takes for the CME front to reach the ParkerSP
  from its onset is computed as: $6.7 + 42.4 = 49.1$ hours.
\item \textbf{True and false statements}
  \begin{description}
  \item[A.] A larger number of selected images would give more precise
    information on the evolution of the CME and the calculated physical
    parameters. TRUE
  \item[B.] The analysis and measurements of the CME evolution should
    consider the differential rotation of the Sun, and therefore the
    calculated velocities will be affected. FALSE
  \item[C.] Any software (numerical) misalignment among the images when
    creating the mosaic will have direct effects on the precision of your
    calculations. TRUE
  \item[D.] The different assumptions made in order to construct the model
    displaying in a heliospheric density map in Figure 1, may affect the
    estimation of the Sun--ParkerSP distance. FALSE
  \item[E.] The interaction of the CME-front with the remnant dust left by
    the 2019 Borisov comet broadens and difuminate the images. This
    intensifies the lack of contrast in the images, increasing the
    uncertainty in determining the CME-front and its propagation. FALSE
    % sic: "difuminate" as in the original
  \end{description}
\item \textbf{Kinetic energy calculations}

  Proton:
  \[
  0.5 \times (1.67262\times10^{-27}~\mathrm{kg}) \times (511~\mathrm{m/s})^2
  = 1.4\times10^{-3}~\mathrm{eV}
  \]
  Alpha particle:
  \[
  0.5 \times (6.64\times10^{-27}~\mathrm{kg}) \times (511~\mathrm{m/s})^2
  = 5.4\times10^{-3}~\mathrm{eV}
  \]
  % sic: the official solution uses 511 m/s (instead of the CME speed of
  % 511 km/s); the quoted eV values follow from the m/s figure.
\end{enumerate}
